\documentclass{report}
\usepackage{amsmath,amssymb}
\setlength{\parindent}{0mm}
\setlength{\parskip}{1em}
\begin{document}
\begin{center}
\rule{6in}{1pt} \
{\large Cameron Musco \\
{\bf Randomized Block Iterative Methods for the Singular Value Decomposition}}

5 Cottage Ct \\ Cambridge MA \\ 02139
\\
{\tt cnmusco@mit.edu}\\
Cameron Musco\\
Christopher Musco\end{center}

In machine learning and data analysis the singular value decomposition is
typically used for principal component analysis, low-rank approximation,
and other problems where just a few top singular directions are needed.
In this regime, block iterative methods like randomized Simultaneous
Iteration have become standard and are often used as defaults in popular
machine learning libraries.

There is good reason for this -- such methods are simple to implement,
fast in practice, and come with strong convergence guarantees for data
analysis tasks. Importantly, unlike traditional bounds, these guarantees
\emph{do not depend on singular values gaps}, which are often helplessly
small in high dimensional data problems.

An analysis by Rokhlin, Szlam, and Tygert of Simultaneous Iteration has
been particularly influential. They demonstrate that, when implemented
with randomly chosen start vectors, the algorithm requires just $\tilde
O(1/\epsilon)$ iterations to return a set of approximate singular vectors
that yield a low-rank approximation within $(1+\epsilon)$ of optimal for
spectral norm error. When requirements on $\epsilon$ are relatively loose
and singular value gaps are small, as is typical in modern data analysis
applications, this iteration bound is much tighter than classical
results.

In this work we extend this gap-independent analysis to a simple
randomized block Krylov method, which is closely related to the classic
Block Lanczos algorithm.
We show that after just $\tilde O(1/\sqrt{\epsilon})$ iterations, our
method recovers a set of approximate singular vectors that are within
$(1+\epsilon)$ of the optimal for low-rank approximation and principal
component analysis. To complement this theoretical guarantee, we
demonstrate experimentally that our method significantly outperforms
randomized Simultaneous Iteration in many scenarios.

Despite their long history, our analysis is the first of a Krylov
subspace method that gives accuracy bounds that do not depend on singular
value gaps. It relies on framing iterative algorithms as de-noising
procedures for coarse \emph{randomized sketching} methods for approximate
singular value decomposition.

Using this framework, we can argue that the approximate singular vectors
returned by our block Krylov algorithm are nearly optimal for low-rank
approximation and principal component analysis even if, due to small
spectral gaps, they have not converged to the true top singular vectors
of the matrix.

Our Krylov method results rely critically on a new analysis of the
standard Rayleigh-Ritz method -- we show that the procedure still works
effectively when convergence to the true top subspace has not occurred.
This new analysis also leads to improved gap-independent bounds for
Simultaneous Iteration.

Finally, beyond gap-independent results, we prove extremely simple
convergence bounds that \emph{do} depend on singular value gaps. We show
how to take advantage of these bounds by using a block size slightly
\emph{larger} than the target rank $k$, theoretically justifying a
heuristic implemented in several data analysis libraries.


\end{document}
